\documentclass[11pt,english]{scrartcl}
\usepackage[T1]{fontenc}
\usepackage[utf8]{inputenc}
\usepackage[margin=2.5cm]{geometry}
\usepackage{babel}
\usepackage{longtable}
\usepackage{array}
\usepackage{xcolor}
\usepackage{listings}
\usepackage[hidelinks,breaklinks=true]{hyperref}

\addtokomafont{disposition}{\rmfamily}
\setlength{\parindent}{0pt}
\setlength{\parskip}{0.5em}
\setlength{\emergencystretch}{3em}
\newcommand{\code}[1]{\nolinkurl{#1}}
\newcommand{\source}[1]{\par{\footnotesize\textbf{Source:} \nolinkurl{#1}}\par}
\makeatletter
\g@addto@macro\UrlBreaks{\do\a\do\b\do\c\do\d\do\e\do\f\do\g\do\h
  \do\i\do\j\do\k\do\l\do\m\do\n\do\o\do\p\do\q\do\r\do\s\do\t
  \do\u\do\v\do\w\do\x\do\y\do\z\do\A\do\B\do\C\do\D\do\E\do\F
  \do\G\do\H\do\I\do\J\do\K\do\L\do\M\do\N\do\O\do\P\do\Q\do\R
  \do\S\do\T\do\U\do\V\do\W\do\X\do\Y\do\Z}
\makeatother
\newcolumntype{L}[1]{>{\raggedright\arraybackslash}p{#1}}
\definecolor{commentgreen}{rgb}{0.1,0.4,0.15}
\lstdefinelanguage{ASL}{
  morekeywords={import,as,conditions,pre,overall,obligation,or,effects,
    always,success,failure,nonperf,infer,observes,onInterrupt,cancel,
    suspend,resume,reset,recovery,actionRecovery,catchingActions,
    postRecovery,retry,continue,goals,excludedGoals,failedActions,
    excludedFailedActions,failureReasons,excludedFailureReasons,
    actionStatuses,maximumAttempts,if,elseif,else,while,for,foreach,
    try,catch,finally,return,exit,async,join,noInterrupt,true,false},
  morecomment=[l]{//},
  morecomment=[s]{/*}{*/},
  morestring=[b]"
}
\lstset{
  basicstyle=\small\ttfamily,
  keywordstyle=\bfseries,
  commentstyle=\color{commentgreen},
  stringstyle=\color{blue!60!black},
  columns=fullflexible,
  keepspaces=true,
  showstringspaces=false,
  breaklines=true,
  breakatwhitespace=false,
  frame=single,
  rulecolor=\color{black!20},
  xleftmargin=0.4em,
  xrightmargin=0.4em,
  captionpos=b
}

\title{DIARC Goal Manager and Action Script Language\\
  Implementation Guide and API Reference}
\author{DIARC contributors\\
  \small Based on the original documentation by Luca Brusatin}
\date{\today}

\begin{document}
\maketitle

\begin{abstract}
The Goal Manager connects symbolic task descriptions to executable DIARC
capabilities. It loads action models, selects actions or plans, manages
concurrent goals, interprets Action Script Language (ASL), checks conditions
and effects, updates belief, and supports interruption, recovery, scheduling,
and action learning. This guide describes the implementation in this
repository rather than a proposed architecture. It separates behavior that
is provided by default from behavior requiring a particular execution
manager, selector, observer, or robot component.
\end{abstract}

\tableofcontents
\clearpage

\section{Scope, terminology, and source map}
\label{sec:architecture}

\subsection{The main concepts}

A \textbf{goal} is a request with an actor, a symbolic predicate, an ID,
priority information, and a lifecycle. An \textbf{action} is a reusable
capability with typed arguments and optional conditions, effects, and
interruption behavior. Actions are either Java primitives exposed through
TRADE or ASL scripts composed from other capabilities. A \textbf{condition}
must hold for execution to be acceptable; an \textbf{effect} describes a
modeled result. An \textbf{observer} tests a predicate and returns variable
bindings. A \textbf{context} is one node in the runtime execution tree,
not necessarily an action: loops, observations, subgoals, and service calls
also have contexts.

The Goal Manager is not itself a navigation, manipulation, perception, or
dialogue implementation. Loading an ASL file that calls a robot capability
does not provide that capability. The corresponding Java/TRADE provider,
belief facts, reference-resolution services, and hardware or simulator
must be supplied by the configuration.

\begin{longtable}{L{0.29\textwidth}L{0.65\textwidth}}
\textbf{Class or subsystem} & \textbf{Responsibility}\\ \hline
\endhead
\code{GoalManagerComponent} &
DIARC component, command-line configuration, public goal services,
notifications, and GUI adapters.\\
\code{manager/ExecutionManager} &
Actor/team hierarchy, pending/active/past goal collections, priorities,
resource arbitration, interpreter threads, learning, and history pruning.\\
\code{goal/Goal}, \code{GoalInfo} &
Goal-form interpretation, actor, identity, status, failure information,
and links to execution contexts.\\
\code{db/Database}, \code{ActionDBEntry} &
Action models loaded from ASL and dynamically discovered TRADE services;
separate action, observer, recovery-policy, and performance databases.\\
\code{selector/ActionSelector} &
Chooses a parameterized action using the selected strategy. Planning is
strategy-dependent, not an automatic property of every submission.\\
\code{ActionInterpreter}, \code{execution/StepExecution} &
Traverses context trees and drives approval, execution, verification,
interruption, and terminal bookkeeping.\\
\code{execution/ActionContext}, \code{GoalContext} &
Action arguments and checks; subgoal selection, goal verification,
and recovery.\\
\code{state/StateMachine} &
Belief-backed state queries and updates, execution history, and branched
states used by non-\code{ACT} executions.\\
\code{asl/ASL.g4}, \code{ActionScriptLanguageParser} &
Language syntax and conversion into \code{ActionDBEntry} roles, metadata,
and \code{EventSpec} instructions.\\
\code{goal/schedule}, \code{manager/ScheduledGoalManager} &
Quartz-backed timed submissions and schedule lifecycle.\\
\end{longtable}

Unless a full repository path is given, Java source paths in this guide
are relative to
\nolinkurl{core/src/main/java/edu/tufts/hrilab/action/}.
ASL resource paths are relative to
\nolinkurl{core/src/main/resources/config/edu/tufts/hrilab/action/asl/}.
Test ASL resources have the corresponding prefix under
\nolinkurl{core/src/test/resources/}.

\subsection{Initialization and runtime organization}

\code{GoalManagerComponent.init()} initializes the state machine and
execution manager before loading the action database. It then loads ASL,
persisted learned actions, optional PDDL and performance models, attaches
notifications, loads priorities, registers the context consultant,
submits command-line goals, and restores configured schedules.

The state machine creates a Belief component using the forwarded belief
arguments. The execution manager obtains the agent/team hierarchy through
\code{getAllAgentTeams}. Consequently, the actor hierarchy must be available
when the execution manager initializes. Merely putting an actor's name in
a goal does not create that actor.

TRADE discovery makes primitives available as providers register.
Execution still requires a live provider; the discovery code registers
for joins and does not provide a matching automatic removal path for
departed primitive providers. Agent-specific capabilities can be
associated with \code{agent:<name>} TRADE groups. The database and action
selector use singleton instances; do not assume that several Goal Manager
components in one JVM have independent action catalogs and selectors.

\source{GoalManagerComponent.java: init, initializeGM, initializeDB}
\source{manager/ExecutionManager.java: init, createInstance}
\source{db/Database.java: addPrimitiveActions, addActionFromTSI}

\section{Getting started and configuring a system}
\label{sec:configuration}

\subsection{Build and launch}

The core uses Java 17 and the repository Gradle wrapper:
\begin{lstlisting}[language=bash]
./gradlew :core:compileJava
\end{lstlisting}

The following small script needs no robot-specific primitives. Save it as
\nolinkurl{core/src/main/resources/config/edu/tufts/hrilab/action/asl/quickstart.asl}:
\begin{lstlisting}[language=ASL,caption={A hardware-independent first script.}]
() = demo["Count three iterations and finish."]() {
  for (!i = 0; !i lt 3; !i++) {
    op:log("info", "Iteration !i");
  }
}
\end{lstlisting}

From the repository root, launch the component with the supplied agent
knowledge base and this script:
\begin{lstlisting}[language=bash]
./gradlew launch \
  -Pmain=edu.tufts.hrilab.action.GoalManagerComponent \
  --args="-beliefinitfile agents/agents.pl -asl quickstart.asl -goal demo(self)"
\end{lstlisting}

The \code{demo(self)} goal explicitly names the actor. Inside ASL,
\code{?actor} is inherited, so the script does not declare or pass an
actor argument to \code{op:log}. The loop prints iterations 0, 1, and 2.
This example demonstrates execution only; an interactive application
normally starts additional components in a Java DIARC configuration.
See, for example,
\nolinkurl{config/src/main/java/edu/tufts/hrilab/config/kinova/KinovaFlexConfig.java}
for loading the Goal Manager alongside capability providers.

\subsection{Action and belief resources}

By default, ASL filenames resolve under the classpath resource directory
\nolinkurl{config/edu/tufts/hrilab/action/asl}. For example,
\code{-asl navigation.asl} loads the resource with that suffix, not an
arbitrary file relative to the shell's working directory. The parser's
file-loading entry point uses \code{getResourceAsStream}.
\code{-dbfileDir} changes the resource-directory prefix.

\code{-asl self:quickstart.asl} restricts actions from that file to
\code{self}; comma-separated actor names before the colon are also
supported. Without a prefix, the file is loaded without an explicit
agent restriction. Restrictions affect selection; they do not create
agents or instantiate robot services.

Belief initialization is separate from ASL loading. At a minimum, supply
appropriate actor/team knowledge. The repository's
\nolinkurl{core/src/main/resources/config/edu/tufts/hrilab/belief/agents/agents.pl}
contains \code{diarcAgent}, \code{memberOf}, \code{object}, and \code{team}
facts/rules used by configurations. It describes many agents and is an
example resource, not a requirement to use that exact hierarchy.

\source{db/Database.java: loadDatabaseFiles, loadDatabaseFromFile}
\source{asl/ActionScriptLanguageParser.java: parseFromFile}

\subsection{Implemented component options}

These are Goal Manager options, in addition to the options inherited from
\code{DiarcComponent}. A filename below means a resource unless persistence
under the user's home directory is explicitly stated.

\begin{longtable}{L{0.36\textwidth}L{0.58\textwidth}}
\textbf{Option} & \textbf{Meaning}\\ \hline
\endhead
\code{-asl <files...>} or \code{--dbfile <files...>} &
Load ASL files, optionally prefixed with \code{agent1,agent2:}.\\
\code{-dbfileDir <dir>} &
Override the default ASL resource prefix.\\
\code{-goal <predicates...>} &
Submit initial goals after initialization. Use
\code{action(actor,args)} or \code{goal(actor,state)} syntax.\\
\code{-selector <class>} &
Fully qualified \code{ActionSelector} subclass. Default:
\code{edu.tufts.hrilab.action.selector.UtilitarianActionSelector}.\\
\code{-executionManagerType <class>} &
Fully qualified execution-manager class. Default:
\code{edu.tufts.hrilab.action.manager.ExecutionManager}.\\
\code{-priorityFile <files...>} &
Priority JSON resources under
\nolinkurl{config/edu/tufts/hrilab/action/manager/priority}.
Default is \code{default.json}. Later files overwrite earlier map entries.\\
\code{-badaction <name>} &
Add an action name to initial forbidden-action constraints.\\
\code{-badstate <predicate>} &
Intended to add a forbidden state. This checkout's \code{parseArgs}
reads the \code{badaction} value in this branch; do not rely on this
option as a working state-safety configuration.\\
\code{-pddl <domain> <problem>} &
Load PDDL actions into the action database, facts into belief, and the
problem goal as an initial goal for \code{self}. The default PDDL
resource prefix is \nolinkurl{config/edu/tufts/hrilab/action/pddl}.\\
\code{-performancefile <files...>} &
Load performance models. Default prefix:
\nolinkurl{config/edu/tufts/hrilab/action/performancemodels}.\\
\code{-performancefileDir <dir>} &
Declared directory override. The current argument parser checks a
lowercase spelling of this option; verify that an override is applied
before relying on it.\\
\code{-learningGui} &
Enable the action-learning GUI flag. This is not the obsolete
\code{-editor} option.\\
\code{-nomemorymanager} &
Disable automatic pruning of old goals, contexts, and state.\\
\code{-historylength <milliseconds>} &
Retention length, default 60000 ms. The pruning cycle is also 60000 ms;
retention is not an exact deletion deadline.\\
\code{-scheduledGoalsFile <filename>} &
Restore/save schedules at
\nolinkurl{~/.diarc/scheduledGoals/<filename>}.\\
\code{-learnedActionsFile <filename>} &
Load and persist taught ASL at
\nolinkurl{~/.diarc/learnedActions/<filename>}.\\
\end{longtable}

The component forwards \code{-beliefinitfile}, \code{-beliefworking},
\code{-beliefepisodic}, \code{-beliefuniversal}, \code{-beliefprover},
\code{-beliefdisk}, \code{-beliefg}, and \code{-beliefgroups} to its local
Belief component. Its own TRADE groups are also added to that component.
The old \code{-agentname}, \code{-beliefcomponent}, \code{-beliefargs},
and \code{-editor} interface described in earlier documentation is not
implemented by this class.

The provided web GUI integration consists of \code{GoalViewerAdapter}
and \code{GoalManagerAdapter}; \code{ActionProgrammerAdapter} and
\code{GoalQueueViewerAdapter} also exist in the action GUI package.
GUI availability depends on the surrounding GUI configuration.

\source{GoalManagerComponent.java: additionalUsageInfo, parseArgs, getAdapterClassNames}

\section{Goal forms and the execution lifecycle}
\label{sec:goals}

\subsection{Action, state, observation, persistent, and scheduled goals}

The outer predicate is significant. The implementation does \emph{not}
guess whether a bare predicate is an effect or an action:

\begin{longtable}{L{0.43\textwidth}L{0.51\textwidth}}
\textbf{Submission predicate} & \textbf{Interpretation}\\ \hline
\endhead
\code{goToLocation(self,kitchen)} &
Explicit action request. Match an action signature, with \code{self} as
the first (actor) argument.\\
\code{goal(self,at(self,kitchen))} &
Desired-state request. Select by effect or use the configured planner
to achieve the inner predicate.\\
\code{obs(self,at(self,kitchen))} &
Observation request. Find an observer for the inner predicate;
\code{observe} is also accepted as the wrapper name.\\
\code{persistent(patrol(self))} &
Repeat the selected execution through a \code{PersistentContext}.
The wrapper takes one complete goal/action predicate.\\
\code{scheduled(patrol(self),period(5,MINUTES))} &
Create a schedule that submits separate executions; see
Section~\ref{sec:scheduling}.\\
\end{longtable}

These navigation/patrol names are illustrative: they must be supplied by
the chosen domain. For an explicit action request the actor is always the
first argument, even if the underlying Java method has no actor parameter.
In \code{goal(actor,state)} the actor is a separate argument from the
desired state, which may itself mention one or more actors.

\code{persistent} means repeated execution, not ``keep a condition true
using an arbitrary maintenance planner.'' In normal acting modes the
persistent context resets its child for another iteration; suspension,
cancellation, and failures still matter. Give repeating behaviors a
blocking operation or an appropriate delay to avoid a tight loop.
The old \code{sti} and \code{did} goal forms are explicitly unsupported.

\source{goal/Goal.java: createGoal, constructors, getGoalPredicate}
\source{execution/PersistentContext.java}

\subsection{What happens after submission}

\begin{enumerate}
\item The component creates a \code{Goal}, attaches its notifier, and
sets execution type, tier, numeric priority, and submitter metadata.
Submission returns an ID; it does not wait for success.
\item The execution manager validates the actor against its hierarchy.
An unknown actor produces a failed goal with an
\code{actorInHierarchy(actor)} justification. A returned positive ID
therefore does not imply that execution started.
\item Scheduled goals go to the scheduler. Ordinary goals may instead be
handed to an active action-learning session. Other goals enter resource
arbitration as pending goals.
\item When allowed by priority and resource policy, the goal becomes
active. Root constraints are checked, a state machine is chosen, and
an \code{ActionInterpreter} is submitted to the executor.
\item A \code{GoalContext} binds the goal and selects an action, or creates
an observation context. Action contexts bind typed arguments and check
constraints, preconditions, and obligations.
\item The interpreter executes primitives or traverses ASL child
contexts. Overall conditions are monitored during progress.
\item Completion verifies return values, effects, and, when the
appropriately typed service is available, norms. State changes are
recorded. State goals additionally verify the desired predicate in the
state machine; explicit action goals do not add a separate desired-state
test beyond the action's own verification.
\item Failures can trigger a catch block or matching recovery policy.
Terminal bookkeeping moves the goal to history, releases resources,
notifies waiters/listeners, and considers other pending goals.
Unexpected interpreter exceptions also fail the goal and run terminal
bookkeeping rather than leaving resources occupied indefinitely.
\end{enumerate}

Subgoals inside an ASL context tree are not the same thing as separately
submitted top-level goals with independently managed public IDs. In
particular, \code{goal:...} composes a subgoal into the current execution;
\code{submitGoal(...)} explicitly asks the execution manager to create
another managed goal.

\source{manager/ExecutionManager.java: submitGoal, executeGoal, goalStatusChanged}
\source{ActionInterpreter.java: call}
\source{execution/GoalContext.java: selectAction, verifyEffects}
\source{execution/StepExecution.java: canDoStep, doStep, finishStep}

\subsection{Goal status versus action status}

\code{GoalStatus} is the public task lifecycle:
\begin{description}
\item[\code{UNKNOWN}] No retained goal with this ID was found.
\item[\code{PENDING}] Waiting to begin or to regain execution resources.
\item[\code{ACTIVE}] Managed execution is active.
\item[\code{SUSPENDED}] Execution was paused. A superseded goal can be in
the pending collection while retaining a suspended execution state.
\item[\code{CANCELED}] Terminal cancellation, not successful completion.
\item[\code{SUCCEEDED}] Terminal successful completion.
\item[\code{FAILED}] Terminal failure; inspect the justification.
\end{description}

Only \code{CANCELED}, \code{SUCCEEDED}, and \code{FAILED} are terminal.
Collection membership and status answer different questions: the active
collection can include explicitly suspended goals, while the scheduled
collection represents schedule controllers.

\code{ActionStatus} is the finer-grained execution/context state. Its
normal phases include \code{INITIALIZED}, \code{CHECKINGAPPROVAL},
\code{APPROVED}, \code{PROGRESS}, \code{VERIFYING_RETURNVALUE},
\code{VERIFYING_EFFECTS}, \code{VERIFYING_NORMS}, and \code{SUCCESS}.
\code{RECOVERY}, \code{SATISFYING_PRECONDITIONS}, and
\code{SATISFYING_OBLIGATIONS} describe recovery/planning-related phases.

Failure statuses include \code{FAIL}, \code{FAIL_FORBIDDEN},
\code{FAIL_CONSTRAINTS}, \code{FAIL_PRECONDITIONS},
\code{FAIL_OVERALLCONDITIONS}, \code{FAIL_OBLIGATIONS},
\code{FAIL_POSTCONDITIONS}, \code{FAIL_ARGUMENTS},
\code{FAIL_ARGUMENTUNBOUND}, \code{FAIL_RETURNVALUE},
\code{FAIL_NOTFOUND}, \code{FAIL_ANCESTOR}, \code{FAIL_CHILD},
\code{FAIL_SYNTAX}, \code{FAIL_RECOVERY}, \code{FAIL_NORMS},
and \code{EXIT_RECOVERY}. Interruption/control states include
\code{CANCEL}, \code{PARENT_CANCEL}, \code{SUSPEND},
\code{PARENT_SUSPEND}, \code{RESUME}, and \code{RETURN};
\code{UNKNOWN} is also defined.

Do not compare the two enums as if their names or granularity matched.
\code{getActionStatus(gid)} reports the goal's root context status, not
necessarily the currently executing leaf primitive.

\source{goal/GoalStatus.java}
\source{ActionStatus.java}

\section{Selection, planning, priorities, and concurrency}
\label{sec:selection}

\subsection{How actions are found}

Candidate lookup uses the goal's kind: observation predicates query the
observer database, explicit action requests query action signatures, and
state goals query effects. Candidates are converted to
\code{ParameterizedAction} instances, checked for valid Java argument
types, and annotated with constraint justifications.

The default \code{UtilitarianActionSelector} sorts candidates using
\code{benefit - cost} and chooses the first constraint-supported candidate,
or a candidate with a failing justification if none is supported. The
latter is rejected during goal execution. \textbf{Implementation caveat:}
the current comparator orders increasing \code{benefit - cost}, despite
the utility-oriented class name. Do not assume it chooses the highest
utility or the fastest action.

The default selector does not solve unsatisfied preconditions, synthesize
arbitrary multistep plans, or guarantee that it will try a second
applicable action after a precondition failure. Preconditions are checked
when the selected action is approved. Use an appropriate planning selector
or explicit recovery if that behavior is required.

\code{PreconditionPlanningSelector} supports planning for a selected
action's preconditions. \code{GoalPlanningActionSelector} first uses that
strategy and then falls back to its \code{FFPlanner} for a goal plan.
The repository also contains random, typed, and MCTS selection strategies.
These are different policies, not interchangeable guarantees.

For FF-based planning, configure the external planner executable using the
repository's \code{diarc.planner.ff} property and provide an adequate
symbolic domain/belief state. Merely loading PDDL or selecting a planner
does not supply missing robot implementations or accurate models of the
physical world. A planning metric supplied with
\code{submitGoalWithMetric} is meaningful only if the selected planning
backend supports that metric.

\source{selector/ActionSelector.java: getCandidateActions}
\source{selector/UtilitarianActionSelector.java}
\source{selector/PreconditionPlanningSelector.java}
\source{selector/GoalPlanningActionSelector.java}

\subsection{Priority assignment and ordering}

Priority is a pair consisting of a \code{PriorityTier} and a \code{long}
value. Higher tiers win before numeric priority is considered. In
ascending order the enum defines:
\begin{lstlisting}
UNINITIALIZED, KIOSK, LOW, NORMAL, HIGH, URGENT, SKIPPENDING
\end{lstlisting}

Within a tier, larger numeric values win; exact ties favor older goals
(smaller IDs). This scheduling priority is separate from an action's
\code{benefit} and \code{cost} selection metadata.

Priority JSON is indexed by the submitted predicate's name and has the
following shape:
\begin{lstlisting}
{
  "default": {"priority": 1, "priorityTier": "NORMAL"},
  "inspect": {"priority": 10, "priorityTier": "HIGH"}
}
\end{lstlisting}

Without an entry, the default is used. If a loaded configuration lacks a
\code{default} entry, the component inserts \code{NORMAL}/1. The supplied
\code{default.json} gives \code{freeze} an urgent priority and puts many
control, dialogue, and introspection operations in \code{SKIPPENDING}.
The current persistent-goal priority helper only unwraps when it sees
more than one argument, although valid persistent goals have one
argument; for deterministic persistent-task priority use an explicit
submission overload rather than expecting inner-name priority lookup.

\source{GoalManagerComponent.java: loadGoalPriorities, getPriorityTierForGoal, getPriorityForGoal}
\source{manager/ExecutionManager.java: compareGoal}

\subsection{Execution-manager choices and resources}

\begin{longtable}{L{0.31\textwidth}L{0.63\textwidth}}
\textbf{Manager} & \textbf{Actual policy}\\ \hline
\endhead
\code{ExecutionManager} &
Default. Ordinary goals generally do not hold an agent-wide execution
resource, so several can execute concurrently. Learning and freeze
resources still constrain execution. A newly submitted blocked,
lower-priority goal is normally failed with an
\code{availableResources(...)} justification; a previously superseded
goal can remain pending for resumption.\\
\code{QueueExecutionManager} &
Uses agent-wide resources for ordinary goals and leaves blocked goals
pending. Disjoint agents can execute concurrently; overlapping
agent/team goals compete. Higher-priority work can supersede lower
priority work. This is a priority/resource queue, not strict FIFO.\\
\code{SingleGoalExecutionManager} &
Adds agent-wide exclusion but inherits the base rejection policy for
new blocked work, rather than the queue manager's wait policy.\\
\end{longtable}

A goal can \emph{require} a resource to be available without holding it,
or \emph{hold} it for exclusive use. Two required-only uses can coexist.
A held resource conflicts with another goal's required or held use.
Pending-goal conflict checks consider higher-priority pending work, not
all pending work indiscriminately.

In the current managers, resources are principally agent/team-wide,
learning, and freeze resources. A team goal accounts for descendants in
the configured hierarchy. This is not a general ASL mechanism for
declaring arbitrary gripper, camera, or room locks. Implement specialized
resource policy in a manager/provider if those distinctions are needed.

A goal can supersede conflicting active goals only if it outranks all of
them. Ordinary superseded executions are suspended, their resources
released, and their interpreter futures retained for later continuation.
Kiosk-tier goals have special cancellation behavior on supersession.
\code{SKIPPENDING} goals bypass the normal resource dependencies in
these managers; reserve that tier for control/background operations that
are safe to run concurrently, not as a workaround for conflicting motion.

\source{manager/ExecutionManager.java: getRequiredResourcesForGoal, getHeldResourcesForGoal, shouldSupersede, supersedeGoals}
\source{manager/QueueExecutionManager.java}
\source{manager/SingleGoalExecutionManager.java}

\section{Goal Manager service reference}
\label{sec:api}

Services in this section are exposed by \code{GoalManagerComponent}
unless otherwise stated. For TRADE, match both the service name and the
exact argument types, especially \code{long.class} versus
\code{Long.class}. A method being a Java method does not by itself make it
a TRADE service, and a TRADE service is not necessarily an \code{@Action}.

\subsection{Submitting work}

\begin{lstlisting}[language=Java]
long submitGoal(Predicate g)
long submitGoal(Predicate g, ExecutionType type)
long submitGoal(Predicate g, ExecutionType type,
                PriorityTier tier)
long submitGoal(Predicate g, ExecutionType type,
                PriorityTier tier, long priority)
long submitGoal(Predicate g, ExecutionType type,
                PriorityTier tier, long priority, Symbol submitter)
long submitGoalWithMetric(Predicate g, Predicate metric)
\end{lstlisting}

The shortest overload uses \code{ACT}; omitted priorities come from the
priority map. In the tier-only overload, \code{null} or
\code{UNINITIALIZED} requests the configured default tier. The fully
specified overload does not repeat that normalization: supply a valid
tier. An omitted submitter is recorded as \code{unknown}.

The metric form stores a predicate such as \code{minimize(totalCost)}
or \code{maximize(reward)} for the planning layer. All forms return a
goal ID, or \code{-1} if the execution manager returns no goal. As
described above, a goal can already have failed by the time its ID is
returned. Retain the ID rather than relying on a later predicate lookup
to identify one particular submission.

\subsection{Waiting without canceling}
\label{sec:waiting}

\begin{lstlisting}[language=Java]
GoalStatus joinOnGoal(long gid)
GoalStatus joinOnGoal(long gid, long millis)
GoalStatus waitForNextGoalStatusUpdate(long gid)
GoalStatus waitForNextGoalStatusUpdate(long gid, long timeout)
GoalStatus waitForNextGoalStatusUpdate(
    long gid, long timeout, GoalStatus startStatus)
void interruptWaitForNextGoalStatusUpdate(long gid)
\end{lstlisting}

\code{joinOnGoal(gid)} waits through pending time and execution until the
managed goal finishes. The timed overload spends its millisecond budget
on both stages and can return \code{PENDING}, \code{ACTIVE}, or another
nonterminal status. An unknown ID returns \code{UNKNOWN}. Use the
no-timeout overload for an indefinite join; do not assume that passing
zero to every timed API has identical semantics.

\code{waitForNextGoalStatusUpdate} waits for a status different from the
starting status, not necessarily a terminal status. Its timeout is in
milliseconds; zero explicitly means indefinite. It returns \code{null}
for an unknown goal, timeout, or interrupted wait, and also for an already
terminal goal when no starting status was supplied. With an explicit
\code{startStatus}, a different current status can be returned
immediately, including a terminal one. This overload avoids a race between
checking a status and beginning the wait, but is not an event history:
rapid intermediate transitions may not be observed.

The interrupt service wakes a status-update wait; it does not cancel the
target goal. The implementation tracks interruption by goal ID, not by
independent waiter identity, so multiple simultaneous waiters for the same
ID should not treat it as a per-caller cancellation channel.

\textbf{A timeout is not goal cancellation.} Inspect the returned status
and explicitly call \code{cancelGoal} if the application requires it.
Joining a recurring schedule waits on the schedule, not merely its first
execution; see Section~\ref{sec:scheduling}.

\subsection{Canceling, suspending, and resuming}

\begin{lstlisting}[language=Java]
boolean cancelGoal(long gid)
void cancelGoals(GoalFilter filter)
boolean suspendGoal(long gid)
boolean resumeGoal(long gid)
List<Long> suspendSystemGoals()
List<Long> suspendSystemGoals(Symbol agentTeam)
List<Long> resumeSystemGoals()
List<Long> resumeSystemGoals(Symbol agentTeam)
List<Long> cancelSystemGoals(Symbol agentTeam)
boolean cancelPendingGoalByIndex(int index)
boolean cancelPendingGoalsByIndices(Predicate indices)
boolean cancelScheduledGoalsByIndices(Predicate indices)
\end{lstlisting}

\code{cancelGoal} searches managed goals and cancels work according to its
current location; cancellation of a schedule also cancels its unfinished
submitted executions. \code{suspendGoal} pauses an existing execution;
a newly pending goal cannot be suspended because it has not started.
\code{resumeGoal} resumes a suspended execution or updates the continuation
state of superseded work. These methods return whether the operation was
accepted, not a promise of instantaneous physical stopping or completion.
Use status observation when synchronization is required.

\code{cancelGoals} is the preferred bulk interface. \code{ACTIVE},
\code{PENDING}, \code{SYSTEM}, and \code{SCHEDULED} select their respective
operations. \code{ALL} and \code{CURRENT} dispatch to the manager's
current-goal cancellation, not deletion of past records.
\code{PAST} and \code{NEXT} are not supported cancellation filters.
Current-goal cancellation also cancels schedules. Active-goal bulk
cancellation preserves actions named \code{listen}; this is a name-based
exclusion, not a blanket exemption for every \code{SKIPPENDING} goal.

System-goal control operations target an unambiguous active task set at an
actor or a single relevant descendant, excluding \code{SKIPPENDING} goals
and schedule controllers. Ordinary child executions created by a schedule
are still ordinary active tasks. These controls are not ``stop every
agent'' operations. If several
descendants have disjoint active tasks, system-goal selection can fail;
use explicit IDs for deterministic multi-agent control.

Pending indices are zero-based in descending priority order; the
multi-index predicate is \code{and(0,2,...)}. Scheduled-goal indices refer
to the scheduler's returned collection, not the pending priority queue.
Indices can change as work completes; ID-based control is preferable.

The no-argument \code{cancelSystemGoals()} and the old
\code{cancelAllCurrentGoals()}, \code{cancelAllActiveGoals()}, and
\code{cancelAllPendingGoals()} remain present but are deprecated for
removal. \code{getCurrentGoalId(Symbol,int)} is also deprecated.

\subsection{Inspecting one goal and its priority}

\begin{lstlisting}[language=Java]
GoalStatus getGoalStatus(long gid)
ActionStatus getActionStatus(long gid)
Justification getGoalFailConditions(long gid)
GoalInfo getGoalInfo(long gid)
Predicate getGoalPredicate(Long gid)
PriorityTier getGoalPriorityTier(long gid)
long getGoalPriority(long gid)
PriorityTier getPriorityTierForGoal(Predicate g)
long getPriorityForGoal(Predicate g)
boolean hasGreaterPriority(long gid, PriorityTier tier, long priority)
Long getGoalId(Predicate pattern)
List<Long> getGoalIds(Predicate pattern)
\end{lstlisting}

\code{GoalInfo} contains \code{gid}, actor, bound/unbound goal fields,
status, start/end timestamps, failure justification, priority/tier,
submitter, metric, persistent/action/observation/learning flags, and
scheduling information when applicable. It is the metadata object
delivered in status notifications.

\code{getGoalPriorityTier} returns \code{UNINITIALIZED} for an unknown
ID; \code{getGoalPriority} returns \code{-1}.
\code{hasGreaterPriority} compares an existing goal against a speculative
new submission, making the existing goal win an exact priority tie.
It returns \code{true} if the ID is unknown, so it is not an existence test.

\code{getGoalId(pattern)} returns \code{-1} if no match exists and one
matching ID otherwise; the choice is not a deterministic ``newest''
selection when several match. Matching uses the submitted goal
predicate's \code{instanceOf} relation, allowing variable patterns.
Use the list form or retain submission IDs when duplicates are possible.

\textbf{Not all queries are null-safe.}
\code{getGoalInfo} and \code{getGoalPredicate} dereference the goal
directly. \code{getActionStatus} returns \code{UNKNOWN} for an unknown ID,
but a known goal without a root context can still make it unsuitable
before execution starts. The unknown-ID failure-conditions query returns
a placeholder justification with \code{reason(unknown)}, not a reliable
failure/existence indicator. Prefer \code{getGoalStatus} or filtered lists
to establish availability, and handle history expiry.

\subsection{Finding collections of goals}

Each of the following families has the same three selection forms:
\begin{lstlisting}[language=Java]
List<Long> getGoalIds(GoalFilter filter)
List<Long> getGoalIds(GoalFilter filter, Predicate pattern)
List<Long> getGoalIdsForActor(GoalFilter filter, Symbol actor)

List<Predicate> getGoalPredicates(GoalFilter filter)
List<Predicate> getGoalPredicates(GoalFilter filter, Predicate pattern)
List<Predicate> getGoalPredicatesForActor(GoalFilter filter, Symbol actor)

List<GoalInfo> getGoalInfos(GoalFilter filter)
List<GoalInfo> getGoalInfos(GoalFilter filter, Predicate pattern)
List<GoalInfo> getGoalInfosForActor(GoalFilter filter, Symbol actor)
\end{lstlisting}

\begin{longtable}{L{0.20\textwidth}L{0.74\textwidth}}
\textbf{Filter} & \textbf{Collection meaning}\\ \hline
\endhead
\code{ACTIVE} & Active collection, including executions explicitly suspended there.\\
\code{PENDING} & Waiting or superseded work.\\
\code{SCHEDULED} & Schedule controllers, not just executions created by them.\\
\code{CURRENT} & Active, pending, and scheduled collections.\\
\code{PAST} & Retained terminal goals.\\
\code{ALL} & Active, pending, past, and scheduled collections.\\
\code{NEXT} & Manager-selected next pending goal.\\
\code{SYSTEM} & Active non-\code{SKIPPENDING} tasks aggregated across the
root actor and its descendants; unlike system-goal controls, this query
does not require a single unambiguous task branch.\\
\end{longtable}

Actor filtering compares actor names while ignoring semantic type; it
does not automatically include all goals of a team's descendants.
Predicate filtering uses \code{instanceOf}. List ordering is not a
general chronological contract. In this checkout the \code{NEXT} filter
can include a null manager result when no next goal exists, so prefer a
\code{PENDING} query for robust empty-queue handling.

\subsection{Preview, introspection, notifications, and local utilities}

\begin{lstlisting}[language=Java]
ParameterizedAction selectActionForGoal(Predicate goalPredicate)
Predicate getActionSteps(Predicate goalPredicate)
void registerGoalManagerNotification(
    TRADEServiceInfo callback, NotificationType type)
void unregisterGoalManagerNotification(
    TRADEServiceInfo callback, NotificationType type)
\end{lstlisting}

\code{selectActionForGoal} invokes the configured selector without
submitting a managed execution. It can return \code{null}.
\code{getActionSteps} uses the selector and returns
\code{steps(step1,step2,...)} for the selected action, or \code{null}.
This describes selected top-level instructions, not a guaranteed,
fully expanded future trace, and does not reserve resources.

Callbacks must be registered TRADE services with these signatures:
\begin{lstlisting}[language=Java]
void onGoalStatus(GoalInfo info)          // GOAL_STATUS_CHANGE
void onActionAdded(ActionDBEntry entry)  // DATABASE_ADDITION
void onActionRemoved(ActionDBEntry entry)// DATABASE_REMOVAL
\end{lstlisting}

Pass the callback's \code{TRADEServiceInfo} and the corresponding
\code{NotificationType}; callback names are application-defined.
Registration is not a request for a historical snapshot. Combine a
snapshot query with notifications when building a UI, and unregister
when the subscriber shuts down. Keep callbacks short: notification
delivery invokes the registered service synchronously in the notifying
thread. \code{getGoalManagerNotifier()} exposes the notifier for internal
integration; ordinary clients should use registration services instead.

The Java method
\code{getActionContextSteps(Predicate goal, Predicate location)}
requires a uniquely matching \code{CURRENT} goal and returns steps at
a location in its context tree. Unlike \code{getActionSteps}, it has no
\code{@TRADEService} annotation in this checkout.
\code{setState(Predicate)} and \code{getGoal(long)} are local testing
utilities, also not public TRADE services.

\source{GoalManagerComponent.java: public service methods, getFilteredGoals}
\source{notification/GoalManagerNotifier.java}
\source{goal/GoalInfo.java}

\section{ASL: action definitions, values, and dispatch}
\label{sec:asl}

\subsection{What ASL compiles into}

ASL is a domain-specific interpreted language, not Java source and not
the AgentSpeak language used by unrelated systems. ANTLR parses a file
into one or more action definitions. The listener builds action-database
roles, metadata, and a linear list of \code{EventSpec} instructions;
runtime control contexts interpret the nested control structure.

Lexical and syntax errors are collected before conversion. A malformed
file is rejected, with parser entry points logging the error and returning
\code{null}; do not expect valid definitions earlier in an invalid file
to be partially loaded.

\subsection{Definition shape, imports, and arguments}

\begin{lstlisting}[language=ASL,caption={Typed inputs, optional input, and an output role.}]
import java.lang.Double;
import java.lang.String;
import edu.tufts.hrilab.fol.Symbol;
import edu.tufts.hrilab.fol.Predicate as Pred;

(Double ?answer) = addOffset["Add an optional offset."](
    Double ?value, Double ?offset = 1.0) {
  ?answer = op:add(?value, ?offset);
}

() = calculationDemo() {
  Double !result;
  !result = act:addOffset(2.0);
  op:log("info", "Result: !result");
}
\end{lstlisting}

The parentheses to the left of \code{=} declare output roles; empty
\code{()} means no outputs. The parentheses after the action name declare
inputs. The optional bracketed string is a human-readable description.
Inputs and outputs use \code{?name}; script-local variables normally use
\code{!name}. A default value makes an input optional. Keep optional
inputs after required inputs and use distinct signatures for meaningful
overloads.

The database inserts a first \code{?actor} role if it is not explicitly
declared. That role has Java type \code{Symbol} and is inherited by ASL
calls by default. It is not a process-wide global variable and does not
require a component \code{-agentname} option. If explicitly declaring
\code{?actor}, it must be the first role.

Imports resolve Java type names; \code{as} provides an alias. They do
not import ASL actions or instantiate components. Actions in any loaded
file are looked up through the database. Use resolvable Java types and
read parser warnings: unresolved type names currently fall back to
\code{String}, which can conceal a type error until dispatch.

\subsection{Variables, bindings, and expressions}

The grammar supports Java type names, generic type syntax, numeric
literals, booleans, quoted strings, symbols, and nested predicates.
However, runtime type resolution uses Java class names directly and does
not erase generic syntax: use raw imported types such as \code{List}
for roles rather than assuming \code{List<Symbol>} resolves to a Java
class. The grammar accepting a type is not proof of runtime support.
Variable names after \code{?} or \code{!} begin with a lowercase letter.
A semantic annotation such as \code{Symbol ?object:physobj} is distinct
from the Java type: \code{Symbol} controls Java dispatch and
\code{physobj} supplies FOL type information.

Declare ordinary locals at the top level of the action body, before
using them:
\begin{lstlisting}[language=ASL]
import java.lang.Double;

(Double ?sum, Double ?difference) = sumAndDifference(
    Double ?left, Double ?right) {
  ?sum = op:add(?left, ?right);
  ?difference = op:subtract(?left, ?right);
}

() = useResults() {
  Double !sum;
  Double !difference;
  (!sum, !difference) = act:sumAndDifference(8.0, 3.0);
  op:log("info", "sum=!sum difference=!difference");
}
\end{lstlisting}

Bindings are shared through \code{ActionBinding} objects; assigning a
callee output can update the caller's destination role. Assign outputs
before normal completion or \code{return;}. There is no
\code{return expression;} construct and no general Java-style statement
\code{!x = !y + 1;}. Use an operator/action/service call on the right-hand
side, for example \code{!x = op:add(!y,1);}. A literal initializer is
allowed in a declaration, such as \code{Double !x = 0.0;}.

Ordinary declarations are not allowed directly inside control blocks by
the \code{controlContent} grammar. Loop variables and catch details have
their own runtime binding rules. A declaration initialized by a call is
translated into a role plus an event; put initialization before its first
use rather than assuming ordinary Java block scope.

Conditions can bind otherwise unbound local or output roles through
their justifications. If several bindings are possible, approval may
choose one arbitrarily with a warning. This is not automatic iteration
over all solutions; use an observer returning a collection and explicit
iteration when all results matter.

The runtime substitutes script variables in strings, as used in the log
examples. For structured arguments prefer actual variables or predicates
over constructing predicate text inside strings. A quoted string and a
\code{Symbol}/\code{Predicate} value are not interchangeable Java types,
even when conversion is available.

\source{asl/ASL.g4: definition, parameter, localvar, specFunction, controlContent}
\source{asl/ActionScriptLanguageParser.java: addRole, enterSpecFunction}
\source{ActionBinding.java}
\source{execution/ActionContext.java: isApproved}

\subsection{The six instruction prefixes}

\begin{longtable}{L{0.25\textwidth}L{0.69\textwidth}}
\textbf{Instruction} & \textbf{Meaning}\\ \hline
\endhead
\code{act:pickUp(?object);} &
Call a named action/script directly. Uses its action model and normal
action verification; does not ask a selector to achieve a state.\\
\code{goal:holding(?actor,?object);} &
Create a desired-state subgoal and invoke goal selection/recovery.
The actor is inherited; \code{?actor} here is part of the desired predicate.\\
\code{obs:visible(?object);} &
Observe a predicate through the observer database. The predicate
arguments are not the observer Java method's explicit parameter list.\\
\code{tsc:someService(?object);} &
Call a TRADE service directly, even if it is not an action.
Does not supply the action model's condition/effect checks.\\
\code{op:add(1,2)} &
Invoke a registered built-in operator, optionally binding its result.
Arithmetic, comparison, collections, FOL, and reflection helpers live
under \code{operators/}.\\
\code{this:getCommand()} &
Invoke a method on the enclosing action context, for introspection.
This is not a method call on a robot component or Java \code{this}.\\
\end{longtable}

Action, goal, observation, and service calls can qualify the actor before
the prefix, for example \code{robotone.act:inspect(?object);} or
\code{?worker.goal:ready(?worker);}. Without qualification the event's
actor is \code{?actor}. Do not add a second actor argument merely because
the external submission form was \code{inspect(robotone,object)}.
For raw \code{tsc} calls, pass exactly the arguments of the Java service;
actor qualification helps service/provider selection and is not an
extra Java argument.

Prefer \code{act} for operations whose modeled success/failure matters.
Use \code{tsc} for deliberate low-level service access. Neither prefix
turns a false boolean data return into a universal exception:
primitive \code{Justification} failures are explicitly verified, while
boolean values can be tested in ASL boolean expressions.

\source{EventSpec.java: EventType, getActor}
\source{execution/Context.java: setupNextStep}
\source{execution/TradeServiceContext.java}
\source{execution/ThisContext.java}

\subsection{Comments and metadata}

Use \code{//} line comments and \code{/* ... */} block comments.
Recognized numeric action settings are \code{benefit}, \code{cost},
\code{minurg}, \code{maxurg}, and \code{timeout}, for example:
\begin{lstlisting}[language=ASL]
() = modeledStep() {
  benefit = 2;
  cost = 1;
  timeout = 5000;
  op:log("info", "Executing modeledStep");
}
\end{lstlisting}

The parser recognizes these names and stores them in the database
entry; arbitrary setting names are not application configuration.
\code{timeout} is timing metadata used by execution/observation paths,
not a guaranteed hard deadline that forcibly interrupts a blocking
hardware service. For bounded external work, implement a timeout and
interruption protocol in the provider and use timed waits appropriately.
Likewise, storing urgency metadata does not change the execution
manager's tier/priority ordering by itself.

\section{Conditions, effects, belief, and observations}
\label{sec:models}

\subsection{Conditions: approval and monitoring}

\begin{lstlisting}[language=ASL,caption={A modeled manipulation action; capability names are illustrative.}]
import edu.tufts.hrilab.fol.Symbol;

() = pickUp(Symbol ?object) {
  conditions: {
    pre infer: gripperOpen(?actor);
    or: {
      pre: on(?object,table);
      pre: on(?object,shelf);
    }
    overall infer: workspaceSafe(?actor);
    obligation infer: permitted(?actor,?object);
  }
  effects: {
    success obs: holding(?actor,?object);
    success infer: not(on(?object,table));
    failure infer: pickupFailed(?actor,?object);
    nonperf infer: pickupCanceled(?actor,?object);
  }
  act:closeGripperOn(?object);
}
\end{lstlisting}

All condition entries must hold (conjunction). An \code{or} block
contains alternatives of a single condition type and requires at least
two entries; one alternative suffices. Separate \code{or} blocks are
conjoined. The grammar accepts \code{pre}, \code{overall}, and
\code{obligation}; it does not currently have a \code{post:} condition
production, even though \code{post} is a reserved token.

\code{pre} conditions are checked during action approval.
\code{obligation} conditions are also checked at approval and fail with
\code{FAIL_OBLIGATIONS} if unmet. They are not a general-purpose automatic
norm-enforcement engine. \code{overall} conditions are monitored while
the action is in progress, with a 250 ms sleep between checks; a false
check produces \code{FAIL_OVERALLCONDITIONS}. This polling interval and
provider latency preclude treating the monitor as a real-time safety
controller.

The example's preconditions do not themselves open a gripper, locate an
object, or grant permission. The selected primitive or a preceding script,
planner, or recovery policy must make the state appropriate.

\subsection{Observation policy: default, forced observation, and inference}

The policy suffix goes between the condition/effect type and the colon.
The following is a model fragment to place inside an action definition:
\begin{lstlisting}[language=ASL]
conditions: {
  pre: visible(?object);
  pre obs: reachable(?object);
  pre infer: authorized(?actor);
}
effects: {
  success: delivered(?object);
  success obs: at(?object,table);
  success infer: recorded(?object);
}
\end{lstlisting}

For a \textbf{condition}, no suffix means: try an observer, and fall
back to \code{StateMachine.holds} if no successful observation execution
is available. A successful observation with a false result is a false
condition; it is not overridden by belief. \code{obs} requires observation
and does not permit a belief-only fallback. \code{infer} queries belief
without requesting an observer. Non-executing simulation paths use state
checks rather than ordinary physical condition observations.

For an \textbf{effect}, no suffix means: try observation, but if no
successful observer execution is available, assume the modeled effect.
\code{infer} also assumes the effect rather than proving it was already
true in belief. \code{obs} requires successful observation of the effect.
This distinction is crucial: a modeled success is not necessarily a
sensor-confirmed physical success. A successfully executed observer that
reports the effect false causes effect verification to fail.

\source{Condition.java: observeCondition}
\source{Effect.java: verify}

\subsection{Effects and state changes}

\code{success} effects apply on action success, \code{failure} effects on
failure, and \code{nonperf} effects on \code{CANCEL}. \code{always} effects
are included independently of the status-specific effect set; they are
not a ``finally'' code block. Effects are interpreted by status-update
paths, so avoid using them to represent irreversible physical operations.

Positive predicates are asserted into the state representation and
\code{not(predicate(...))} represents a negative/retraction update.
These are symbolic changes, not calls to actuators. Observations update
belief through observation contexts; action state updates omit effects
already handled by those observers.

Verified success effects and the resulting state support desired-state
goals and planning. Internal generated \code{succeeded(...)} and
\code{failed(...)} predicates also record execution outcomes; application
state goals should normally refer to domain predicates such as
\code{holding} or \code{at}, not depend on generated bookkeeping shapes.

When an observed postcondition fails, other unobserved modeled success
effects can still be recorded by the action context. Neither
\code{FAILED} nor \code{CANCELED} means ``the world was rolled back.''
Recovery must account for partial physical and symbolic changes.

At completion, if TRADE exposes
\code{getViolations(Symbol, StateMachine)}, the step executor checks its
\code{Justification}; violations and service-call errors fail the step
instead of silently declaring norm success. Without that service,
completion does not perform that norm check.

\source{db/ActionDBEntry.java: getEffects}
\source{execution/ActionContext.java: getStateUpdates, verifyEffects}
\source{execution/StepExecution.java: finishStep}
\source{state/StateMachine.java}

\subsection{Implementing an observer}

An observer is a TRADE service annotated with \code{@Observes}; it takes
one \code{Term} or \code{Predicate} query and returns
\code{List<Map<Variable,Symbol>>}. The query is a predicate to evaluate,
not simply an object name:
\begin{lstlisting}[language=Java]
@TRADEService
@Observes({"at(?object,?place)"})
public List<Map<Variable, Symbol>> observeAt(Term query) {
  // Query the configured sensor or world model.
  // Return bindings for every solution matching query.
  return sensorBindings(query);
}
\end{lstlisting}

Here \code{sensorBindings} is application-specific pseudocode, not a
repository helper. A nonempty list means the positive predicate was
observed and supplies its bindings. For a true ground query, return a
list containing one empty map. An empty list means no positive solution;
\code{null} means observation execution failed or could not be completed.
These are not equivalent. Negated queries are evaluated using the
positive observer results.

An ASL observer can declare \code{observes: predicate(...);} and must
provide compatible query/result roles. For working patterns consult
\nolinkurl{core/src/test/resources/config/edu/tufts/hrilab/action/asl/observation_tests.asl}.
The grammar accepts one predicate per \code{observes} declaration.

\source{annotations/Observes.java}
\source{execution/ObservationContext.java: verifyReturnValue, getBoundObservationResults}

\section{ASL control flow and concurrency}
\label{sec:control}

\subsection{Boolean expressions and branching}

\begin{lstlisting}[language=ASL]
() = conditionalDemo() {
  if (op:gt(4,2) && ~op:equals("a","b")) {
    op:log("info", "First branch");
  } elseif (false) {
    exit(FAIL);
  } else {
    op:log("info", "Fallback branch");
  }
}
\end{lstlisting}

Boolean expressions accept \code{true}, \code{false}, prefixed calls,
parentheses, \code{~}, \code{&&}, and \code{||}. Conjunction binds more
tightly than disjunction, and the control contexts short-circuit.
The keyword is \code{elseif}, not Java's two-word \code{else if}.
Arithmetic and comparisons are calls such as \code{op:lt(!x,10)};
ordinary infix Java expressions are not accepted as arbitrary conditions.
An action's boolean return can decide a branch even though a false boolean
return is not generally a primitive-action failure outside a boolean test.

\subsection{Loops}

\begin{lstlisting}[language=ASL]
import java.lang.Integer;
import java.util.List;

() = loopDemo() {
  Integer !remaining = 3;
  while (op:gt(!remaining,0)) {
    op:log("info", "remaining=!remaining");
    !remaining = op:decrement(!remaining);
  }
  for (!i = 0; !i lt 3; !i++) {
    op:log("info", "index=!i");
  }
}

() = visitItems(List ?items) {
  foreach (!item : ?items) {
    op:log("info", "item=!item");
  }
}
\end{lstlisting}

\code{for} is an integer loop, not a general Java \code{for} clause.
The initialization, comparison, and update must use the same loop
variable; the parser rejects a mismatch. Supported runtime comparisons
include \code{lt}, \code{leq}, \code{gt}, and \code{geq}; updates include
\code{++}, \code{--}, \code{+=}, \code{-=}, \code{*=}, and \code{/=}.
The initial value can be omitted for an existing loop variable.

\code{foreach} expects a variable bound to a Java \code{Collection}.
It binds one element per iteration; order follows the supplied
collection, so a set does not imply deterministic order.
The grammar has no \code{break} or \code{continue} statements.
\code{return;} exits the enclosing action rather than only the loop.
Loops discard previous iteration child contexts, so the context tree is
not a complete per-iteration event log.

\subsection{Return, explicit failure, and structured handling}

\code{return;} completes the current script early, using already bound
output roles. \code{exit(FAIL);} explicitly fails; a more specific
\code{ActionStatus} and predicates/justifications may be supplied for
diagnostics, for example \code{exit(FAIL_PRECONDITIONS,ready(?actor));}.
These arguments are used to construct failure justifications, not simply
printed as error strings.

\begin{lstlisting}[language=ASL,caption={Catch an expected action status and run cleanup.}]
() = handlingDemo() {
  try {
    exit(FAIL_NOTFOUND);
  } catch(FAIL_NOTFOUND, !details) {
    op:log("warn", "Handled missing capability: !details");
  } finally {
    op:log("info", "Cleanup");
  }
}
\end{lstlisting}

\code{catch(STATUS)} handles that status; \code{catch()} is a catch-all.
\code{catch(!details)} or \code{catch(STATUS,!details)} binds a local
\code{Justification}. This is status-based ASL handling, not Java's
exception-type syntax. An unmatched failure propagates.
\code{finally} runs for normal completion and failure, and the repository
also tests return-through-finally behavior. Cleanup can itself fail.
Do not swallow a failure unless the handler actually repairs it or the
failure is intentionally acceptable.

\subsection{Asynchronous blocks and joins}

\begin{lstlisting}[language=ASL]
import java.lang.Long;
import edu.tufts.hrilab.action.ActionStatus;

() = parallelDemo() {
  Long !asyncId;
  ActionStatus !status;
  ActionStatus !expected = SUCCESS;
  !asyncId = async {
    op:sleep(100);
    op:log("info", "Background work complete");
  }
  op:log("info", "Foreground work");
  !status = join(!asyncId);
  if (~op:equals(!status,!expected)) {
    exit(FAIL);
  }
}
\end{lstlisting}

\code{async} starts a child context on another interpreter thread and can
return its \code{Long} context ID. \code{join(!asyncId)} waits for that
context; \code{join(!asyncId,timeout)} uses a millisecond timeout.
The optional return value is \code{ActionStatus}, potentially nonterminal
after a timeout. Check it explicitly when the result matters.

An ASL async ID is \emph{not} a managed goal ID:
\code{join(...)} searches the enclosing action's context tree, whereas
\code{joinOnGoal(...)} joins an execution-manager goal. Do not interchange
the IDs. Concurrent child contexts also do not independently acquire
top-level agent resources. Avoid concurrent calls to the same hardware or
unsynchronized mutation of shared bindings; joining is not mutual exclusion.

\subsection{Interruption hooks and protected-looking blocks}

\begin{lstlisting}[language=ASL,caption={Illustrative provider-specific interruption protocol.}]
import edu.tufts.hrilab.fol.Symbol;

() = driveTo(Symbol ?destination) {
  onInterrupt: {
    cancel: tsc:stopMotion();
    suspend: tsc:pauseMotion();
    resume: reset;
  }
  act:performDrive(?destination);
}
\end{lstlisting}

The application must provide the named services/actions. A hook specifies
one \code{act}, \code{goal}, or \code{tsc} call per event, optionally
actor-qualified. \code{resume: reset;} requests reinitialization of the
action execution rather than assuming an interrupted primitive can resume
where it stopped. Design retries to account for already completed side
effects.

The grammar also recognizes a \code{noInterrupt} block and creates a
\code{NoInterruptContext}. In this checkout that context is essentially
a sequential block: it does not implement a general cancellation mask or
resource lock. Do not use its name as evidence that hardware operations
cannot be interrupted, or that concurrent goals are excluded.
Actual stopping/resuming remains a provider/interpreter integration
responsibility.

\source{execution/control/IfContext.java, AndContext.java, OrContext.java}
\source{execution/control/ForContext.java, ForeachContext.java}
\source{execution/control/TryContext.java, CatchContext.java, FinallyContext.java}
\source{execution/control/AsynchronousContext.java, JoinContext.java, NoInterruptContext.java}
\source{asl/ActionScriptLanguageParser.java: enterOnCancel, enterOnSuspend, enterOnResume}

\section{Writing Java primitives}
\label{sec:primitives}

Register a service with TRADE and mark it with \code{@Action} to make it
an ordinary action primitive. The example below is a provider template;
the motion driver is application-specific:
\begin{lstlisting}[language=Java]
@TRADEService
@Action
@Condition(
    condition = {"connected(?actor)"},
    type = ConditionType.PRE,
    inferable = {"connected(?actor)"})
@Effect(
    effect = {"at(?actor,?destination)"},
    type = EffectType.SUCCESS,
    observable = {"at(?actor,?destination)"})
@OnInterrupt(
    onCancelServiceCall = "stopMotion()",
    onSuspendServiceCall = "pauseMotion()")
public Justification navigate(Symbol destination) {
  return motionDriver.navigate(destination);
}
\end{lstlisting}

The provider must be instantiated and registered, not merely on the
classpath. The repository compiles with \code{-parameters}; role names
are derived from Java parameter names, so keep annotation variables in
sync with those names. The database inserts an actor role if the Java
method does not have one; when the Java method's first parameter is named
\code{actor}, that explicit actor is passed through instead.

Use the annotations from \code{edu.tufts.hrilab.action.annotations} and
\code{ai.thinkingrobots.trade.TRADEService}. The relevant enum types are
\code{ConditionType} and \code{EffectType}. Place annotations where TRADE
can discover them on the registered provider/interface.

\begin{description}
\item[\code{@Action}] Makes the service an action primitive.
\item[\code{@Condition}] Repeatable. Each annotation is conjoined with the
others; its \code{condition} array is a disjunction. \code{observable}
forces observation of listed strings, and \code{inferable} forces belief
checks. The strings must match the condition entries.
\item[\code{@Effect}] Repeatable. Its \code{effect} array describes
status-specific state changes. \code{observable} forces observation;
the inference-only array is named \code{infer}, not \code{inferable}.
The old \code{sideEffect} attribute is not in this annotation.
\item[\code{@Observes}] Declares observer predicate patterns; follow the
query/result contract in Section~\ref{sec:models}.
\item[\code{@OnInterrupt}] Names service calls for cancellation, suspension,
and resumption. The methods must actually unblock/stop the provider as
appropriate; changing a goal status is not a physical stop command.
\end{description}

For execution failure with an explanation, return a false
\code{Justification}; \code{ActionContext.verifyReturnValue} propagates
it as a failure. A false \code{boolean} is often useful as data or a
branch test, but is not generally verified as a failed primitive action
at completion. A \code{void} primitive with no conditions or observed
effects can appear successful once the service returns; choose a richer
contract for fallible operations.

Although an \code{@RecoveryPolicy} annotation type exists, the primitive
annotation-import path in \code{Database.addActionFromTSI} handles
conditions, effects, observations, and interrupts, not that annotation.
Use the implemented ASL recovery-policy mechanism below rather than
assuming the annotation automatically installs a policy.

\source{db/Database.java: addActionFromTSI, addCondition, addEffect}
\source{db/ActionDBEntry.java: executeAction}
\source{execution/ActionContext.java: verifyReturnValue}
\source{annotations/Action.java, Condition.java, Effect.java, Observes.java, OnInterrupt.java}

\section{Recovery: repairing failure rather than hiding it}
\label{sec:recovery}

Recovery is separate from a local \code{try}/\code{catch}. Runtime
contexts consult the \code{Recovery} machinery, match registered policies
against failure information, and execute a policy under a recovery
context. A recovery policy is an action definition with a \code{recovery}
block; the database stores it in the recovery-policy database rather than
as an ordinary selectable domain action.

\begin{lstlisting}[language=ASL,caption={Action recovery template; assumes openGripper is available.}]
import java.util.List;
import edu.tufts.hrilab.fol.Predicate;

() = repairPickup(List ?failedActionPredicates,
                  List ?failureReasons,
                  Predicate ?failedGoal) {
  recovery: {
    actionRecovery: {
      postRecovery: retry
    }
    failedActions: {pickUp(X,Y)}
    actionStatuses: {FAIL_PRECONDITIONS}
    maximumAttempts: 2
  }
  act:openGripper();
}
\end{lstlisting}

This template matches the \code{pickUp(actor,object)} signature. Capital
\code{X} and \code{Y} are FOL pattern variables in recovery matching,
not automatically bound ASL \code{?} roles. Policy execution receives
the failed-action predicates, failure reasons, and failed goal through
the three input roles shown. Supply domain-specific checks if opening a
gripper is not an appropriate repair for every matched precondition
failure.

\subsection{Policy matching and attempt limits}

The following fields constrain applicability:
\begin{description}
\item[\code{goals} / \code{excludedGoals}] Goal patterns to include or exclude.
\item[\code{failedActions} / \code{excludedFailedActions}] Failed action
signature patterns to include or exclude.
\item[\code{failureReasons} / \code{excludedFailureReasons}] Predicate
patterns describing acceptable/unacceptable failure reasons. A policy
must account for all reasons of a multi-reason failure.
\item[\code{actionStatuses}] An optional set of \code{ActionStatus} names.
\item[\code{maximumAttempts}] Positive integer, default 1, limiting repeated
use for the same failure reasons under the same context.
\end{description}

An include field and its exclude counterpart must not be specified
together; the parser warns and ignores the exclusion when the inclusion
is present. Empty constraints are broad, not a safe default for an
arbitrary robot recovery. Narrow policies to failures they can repair.

Recovery metadata uses predicates in braces separated by commas and does
not use ordinary statement semicolons after fields. This differs from
\code{conditions} and \code{effects} syntax.

\subsection{Goal recovery versus action recovery}

Without an \code{actionRecovery} block, a policy is a goal-recovery policy:
after repair, the goal context can select/plan again and must still
satisfy a desired-state goal.

An \code{actionRecovery} block selects action-level recovery and requires
\code{postRecovery: retry} or \code{postRecovery: continue}.
\code{retry} resets and reattempts the caught action. \code{continue}
treats recovery as having handled that action and permits execution to
continue. Optional \code{catchingActions} patterns in braces delay application
until failure reaches a matching enclosing action, rather than always
repairing the nearest failed call.

For action retry, attempt accounting can reset after the retried action
succeeds; goal-recovery attempt counts do not have that success reset.
Policies may also return \code{PostRecoveryExecutionType} to override the
configured action-recovery continuation. See the executable examples in
\nolinkurl{core/src/test/resources/config/edu/tufts/hrilab/action/asl/recovery/},
particularly \code{actionRecoveryTests.asl},
\code{goalRecoveryTests.asl}, and \code{shouldAttemptRecoveryTest.asl}.

A successful recovery policy is not proof that the original physical
task succeeded. Choose \code{continue} only when it is semantically
correct to skip retrying the caught action. Repeated retries of a
non-idempotent action can duplicate deliveries, motion, or assertions.

\source{recovery/Recovery.java}
\source{recovery/PolicyConstraints.java}
\source{db/RecoveryPolicyDatabase.java}
\source{execution/GoalContext.java, ActionContext.java}

\section{Timed, periodic, and persistent behavior}
\label{sec:scheduling}

\subsection{Schedule predicates}

Submit \code{scheduled(innerGoal,timing...)} using the normal
\code{submitGoal} service. The inner predicate is a complete action or
goal request with its actor:
\begin{lstlisting}
scheduled(inspect(self),start(10,SECONDS))
scheduled(inspect(self),period(5,MINUTES))
scheduled(inspect(self),start(10,SECONDS),
          period(5,MINUTES),end(1,HOURS))
\end{lstlisting}

Provide a start or a period (or both). With only a period, scheduling
starts immediately. An end requires a period. Numeric time-unit forms
use \code{java.util.concurrent.TimeUnit} names, case-insensitively:
\code{MILLISECONDS}, \code{SECONDS}, \code{MINUTES}, \code{HOURS}, etc.
Use plural enum names such as \code{MINUTES}, not natural-language
\code{minute}. Numeric inputs are converted to integral counts before
conversion to milliseconds: use \code{period(1500,MILLISECONDS)} rather
than \code{period(1.5,SECONDS)} for a 1.5-second interval.

Start and end can also use date strings, including
\code{MM-dd-yyyy HH-mm-ss}; consult \code{ScheduledGoal.parseDateString}
for supported alternatives and local-time behavior. A period is a
numeric amount plus a unit, not an arbitrary date string. Validate
inputs before submission; malformed forms may throw during construction
or leave invalid timing state, rather than reliably returning a failed
managed goal.

\subsection{Controller IDs and execution IDs}

A schedule is a long-lived controller with its own ID and a Quartz job.
Each trigger creates a new ordinary goal with a different ID. If the
previous execution has not terminated, the trigger does not submit
another overlapping execution of that schedule. This is not an unbounded
backlog of missed ticks.

Periodic child executions are assigned \code{LOW} tier unless the
configured priority helper returns \code{SKIPPENDING}; their numeric
priority is looked up separately. The controller's explicitly supplied
priority is therefore not a general override for every child execution.
Do not assume other submission metadata, such as a simulation execution
type, is automatically inherited by scheduled children.

When restoring a periodic schedule whose start is in the past, the
scheduler advances to a matching future interval using the full
millisecond period. It does not truncate intervals to whole seconds.

\subsection{Schedule-control services and persistence}

\begin{lstlisting}[language=Java]
boolean updateScheduledGoal(long goalId, Predicate update)
Date getNextScheduledExecutionDate(long goalId)
\end{lstlisting}

Updates accept \code{start(...)}, \code{period(...)}, \code{end(...)}, or
a conjunction, for example
\code{and(start(30,SECONDS),period(10,MINUTES))}. The service returns
\code{false} for a missing schedule or unsuccessful update/reschedule.
Updates are not transactional: a conjunction can update earlier fields
before a later field fails. The next-date query returns \code{null} for
a non-scheduled ID or a schedule with no future firing.

\code{cancelGoal(controllerId)} deletes the schedule and cancels its
unfinished submitted goals. Canceling only one child ID does not remove
the recurring schedule. Joining the controller waits for schedule
completion/deletion, not merely one child; use a timeout for an
unbounded recurring schedule. Scheduled joins retain deletion
notifications even if deletion preceded the wait.

With \code{-scheduledGoalsFile}, schedule predicates are restored at
startup and written through the scheduler's persistence path at shutdown.
This file is a Java-serialized collection of predicates, not an ASL
script or a text calendar. It persists timing/task descriptions, not
live interpreter threads or stable execution IDs across restarts.

For continuous patrol use a persistent action with explicit delay;
for periodic inspection use a schedule. Combining persistence and a
periodic schedule often prevents later ticks from running because the
first child never terminates.

\source{goal/schedule/ScheduledGoal.java}
\source{manager/ScheduledGoalManager.java}
\source{goal/schedule/JoinOnGoalListener.java}

\section{Action learning, freezing, simulation, and history}
\label{sec:advanced}

\subsection{Teaching and editing actions}

Action learning is managed per \code{AgentTeam}. During an active
learning session, ordinary submitted goals for that actor or an
appropriate ancestor can be recorded as taught steps rather than
executed as independent requests. Learning resources keep incompatible
work from silently changing the teaching state; overlapping
ancestor/descendant learning sessions are restricted.

The following services are registered by \code{ExecutionManager}:
\begin{lstlisting}[language=Java]
boolean learnAction(Symbol actor, Predicate newAction)
void waitForActionLearningStart(Symbol actor, Predicate newAction)
boolean pauseActionLearning(Symbol actor, Predicate newAction)
boolean resumeActionLearning(Symbol actor, Predicate newAction)
boolean endActionLearning(Symbol actor, Predicate newAction)
boolean cancelActionLearning(Symbol actor, Predicate newAction)
boolean pauseActionLearning(Symbol actor)
boolean resumeActionLearning(Symbol actor)
boolean endActionLearning(Symbol actor)
boolean cancelActionLearning(Symbol actor)
ActionLearningStatus getLearningStatus(Symbol actor)
void changeLearningExecution(Symbol actor, Symbol status)
void modifyAction(Symbol actor, Predicate action,
                  Predicate modification, Predicate location)
boolean waitForLearningStateUpdate(Symbol actor, Predicate action)
void interruptWaitForLearningStateUpdate(Symbol actor, Predicate action)
Justification forgetTaughtAction(Predicate actionSignature)
List<ActionDBEntry> getLearnedActions()
boolean isLearnedAction(Predicate actionSignature)
\end{lstlisting}

\code{learnAction} is a blocking teaching operation; ending/canceling
the session must be possible from another execution/control path.
The \code{waitForActionLearningStart} service allows the teacher to
synchronize before sending steps. Pause stops accepting steps into the
active teaching state; resume continues it. End builds the taught
action; cancel abandons the current learning operation. The actor-only
overloads use that actor's active teaching signature.

\code{changeLearningExecution(actor,execute)} enables execution while
learning; other status names disable it in the current implementation.
\code{modifyAction} edits a learning state's conditions, effects, or
steps using predicate operations. Consult
\code{learning/ActionLearning.java} for accepted edit/location shapes
and unsupported argument edits rather than treating it as arbitrary
ASL source replacement.

The \code{getLearnedActions} service reads the configured persisted
definitions, rather than returning every action in the live database.
\code{isLearnedAction} checks matching database entries against the taught
action/file criteria. The Java helper \code{loadLearnedActions(boolean)}
is not a TRADE service. With \code{-learnedActionsFile}, taught
definitions are loaded and persisted as ASL. \code{forgetTaughtAction}
is intended for this persisted taught-action workflow and returns a
justification; it is not a generic deletion mechanism for primitives.
Teaching builds action structure from supplied steps, not an automatic
learned motion controller or a complete physical effect model.

\source{manager/ExecutionManager.java: action-learning services}
\source{learning/ActionLearning.java}
\source{learning/LearningState.java}

\subsection{Freeze and unfreeze}

\code{freeze(Symbol actor)} and \code{endFreeze(Symbol actor)} are
execution-manager action services. Freeze blocks through an actor/team
freeze mechanism and interacts with the freeze resource; unfreeze
signals it. In the supplied priorities, freeze is urgent and endFreeze
skips normal pending/resource gating so the frozen system can be
released. This mechanism is for coordinated task control, not a hardware
emergency stop.

\subsection{Execution types}

\begin{longtable}{L{0.30\textwidth}L{0.64\textwidth}}
\textbf{Type} & \textbf{Intended execution path}\\ \hline
\endhead
\code{ACT} & Execute primitives and update the main state machine.\\
\code{OBSERVE} & Nominal observation path for modeled actors, subject to
the actor-based conversion described below. Not the same as submitting
one \code{obs(actor,predicate)} query.\\
\code{SIMULATE} & Use branched symbolic state without ordinary physical
primitive execution.\\
\code{SIMULATE_ACT} & Acting simulation path; primitives can execute.
Requires correctly routed simulation services/environment.\\
\code{SIMULATE_PERFORMANCE} & Use performance models to sample effects
and timing for assessment rather than ordinary execution.\\
\end{longtable}

The execution manager branches state for any submission whose type is
not \code{ACT}. This does not make arbitrary provider code side-effect
free, nor does it fully simulate sensors, Java services, or physical
dynamics. In particular \code{SIMULATE_ACT} is explicitly allowed to
execute calls. Validate provider routing before using it with real
hardware.

Action contexts adapt execution based on actor membership:
\code{setExecType} switches non-DIARC actors to \code{OBSERVE}, and
switches \code{OBSERVE} back to \code{ACT} for DIARC actors. Consequently,
submitting a DIARC actor with \code{OBSERVE} is \emph{not} a reliable
``never actuate'' mode. Actor identity and DIARC-agent membership affect
execution semantics as well as dispatch; use an explicit observation
goal when the intent is to query a predicate.

\subsection{History and shutdown}

Past goals, terminated context trees, and state history are pruned on the
manager's periodic memory task when enabled. The default retention window
is 60000 ms; configure \code{-historylength} or
\code{-nomemorymanager} for longer introspection sessions, with the
corresponding memory cost. A previously valid goal ID can become
\code{UNKNOWN} after pruning; task IDs are not durable audit records.

Component shutdown cancels pending and active goals, waits for active
interpreters, deregisters the context consultant and execution-manager
services, shuts down manager infrastructure, and destroys the database
and selector singletons. A provider that never responds to cancellation
can delay shutdown; implement interruption protocols for blocking work.

\source{execution/ExecutionType.java}
\source{execution/ActionContext.java: setExecType}
\source{manager/ExecutionManager.java: freeze, endFreeze, pruneOldData, shutdown}
\source{GoalManagerComponent.java: shutdownComponent}

\section{Worked integration scenarios}
\label{sec:usecases}

\subsection{A Java client that submits and checks a task}

The following client fragment uses the hardware-independent
\code{demo} script from Section~\ref{sec:configuration}; it assumes
TRADE and the Goal Manager have already initialized:
\begin{lstlisting}[language=Java]
Predicate request = Factory.createPredicate("demo(self)");
long gid = TRADE.getAvailableService(new TRADEServiceConstraints()
    .name("submitGoal").argTypes(Predicate.class))
    .call(Long.class, request);
if (gid < 0) {
  throw new IllegalStateException("Goal submission rejected");
}

GoalStatus status = TRADE.getAvailableService(
    new TRADEServiceConstraints().name("joinOnGoal")
        .argTypes(long.class, long.class))
    .call(GoalStatus.class, gid, 5000L);

if (status == GoalStatus.FAILED) {
  Justification reason = TRADE.getAvailableService(
      new TRADEServiceConstraints().name("getGoalFailConditions")
          .argTypes(long.class))
      .call(Justification.class, gid);
  throw new IllegalStateException("Goal failed: " + reason);
}
if (status != GoalStatus.SUCCEEDED) {
  throw new IllegalStateException("Goal did not succeed: " + status);
}
\end{lstlisting}

This fragment propagates \code{TRADEException} to its caller. In an
application, handle it through the component's logging/error protocol.
A nonterminal result after 5000 ms is not a failed goal; the last branch
reports that this client did not obtain successful completion within its
wait. It deliberately does not cancel the still-running task.
To cancel on timeout, explicitly call \code{cancelGoal(long)} and then
observe the terminal status.

When multiple managers or providers advertise the same service, add
appropriate TRADE group constraints rather than relying on an arbitrary
match.

\subsection{Navigation: command versus requested outcome}

Suppose a provider offers \code{navigate(Symbol destination)} with an
effect \code{at(?actor,?destination)}. A client submits
\code{navigate(self,kitchen)} to request that particular capability, or
\code{goal(self,at(self,kitchen))} to request the outcome. A navigation
observer can confirm the latter predicate from localization.

The outcome form allows the configured selector/planner to choose an
appropriate modeled method. It does not by itself guarantee shortest
paths, automatic obstacle avoidance, or precondition repair. The
navigation subsystem handles motion; the Goal Manager handles task-level
selection and verification. Use observed effects when the physical
arrival must be confirmed instead of assumed.

\subsection{Pick-and-place: compose goals and recover locally}

For a manipulation domain, a script can request
\code{goal:holding(?actor,?object)} and then
\code{goal:on(?object,?destination)} rather than hard-coding every motor
step. Domain actions must advertise corresponding effects, an observer
must confirm the required physical results, and belief must contain
the relevant symbolic objects and relations.

An expected, local sensor failure can be handled in an ASL catch block.
A reusable repair such as opening a gripper before retrying pickup
belongs in a narrowly matched recovery policy. Do not let a catch-all
turn a failed grasp into successful delivery. Repository domain examples
are in \code{demos/pickAndPlace/} and \code{manipulation.asl}; their
capabilities require the associated robot/perception configuration.

\subsection{Multiple robots and operator interruption}

With \code{QueueExecutionManager}, submit independent tasks for
\code{robotone} and \code{robottwo} to allow concurrent work on disjoint
agent resources. A task for their parent team can conflict with both.
An urgent task can supersede lower-priority conflicts only if it outranks
every affected goal; the suspended tasks can later resume.

Use \code{robotone.act:...} and \code{robottwo.act:...} for actor-qualified
steps inside a team script, and implement provider interruption hooks.
Inside-script \code{async} adds context concurrency, not additional
top-level resource arbitration. For operator controls, keep cancel,
resume, and unfreeze paths runnable while the main task is blocked.
Consult \code{demos/multiagent.asl}, \code{demos/queueEM.asl}, and
\code{nao/dance_multiagent.asl} as domain examples, not universal hardware
independence.

\subsection{Recurring inspection and background listening}

Use \code{scheduled(inspect(self),period(5,MINUTES))} for a periodic
inspection that should not overlap its own previous execution.
Keep the controller ID for changing or canceling the recurring task;
track child execution IDs to report individual failures.

For a long-running listener or patrol loop, a persistent goal can repeat
a blocking/delayed action. The provided priority map gives \code{listen}
\code{SKIPPENDING} status so dialogue can coexist with task execution.
Do not copy this bypass priority to arbitrary motion tasks. Dialogue
scripts require their speech/language providers; \code{core.asl} alone
is not a complete spoken-interaction configuration.

\subsection{Teaching a reusable procedure}

Start \code{learnAction(actor,newAction)} through a separate control
execution, wait until learning starts, and submit the intended steps for
that actor. Choose whether the steps should also execute during teaching.
End the session to construct the action, inspect its steps/model, and
use the learned-actions persistence option if the procedure should survive
restart. Learning a step sequence does not automatically infer all
conditions, effects, observability requirements, or safe retry behavior.

\section{Diagnostics, validation, and maintenance}
\label{sec:diagnostics}

\subsection{Common symptoms and concrete checks}

\begin{longtable}{L{0.30\textwidth}L{0.64\textwidth}}
\textbf{Symptom} & \textbf{Check}\\ \hline
\endhead
No action found &
Confirm provider registration, \code{@Action}, ASL resource loading,
actor restrictions/groups, action signature, and Java argument types.
Distinguish \code{act} from \code{goal}.\\
Unknown actor failure &
Load the agent/team belief facts before manager initialization; check
\code{getAllAgentTeams} and the requested actor name.\\
Precondition failure despite a loaded action &
Inspect the justification and its bindings. Loading an action does not
make its preconditions true; the default selector does not plan them.\\
Postcondition failure &
Check observer availability, the query/result contract, and \code{obs}
versus inference policy. Verify that the effect represents the actual
physical result.\\
Success without physical confirmation &
Look for unobserved or \code{infer} effects, \code{void}/boolean primitive
contracts, and raw \code{tsc} calls bypassing an action model.\\
Goal stays pending &
Inspect priority, held resources, team overlap, freeze, learning, and
superseded work. Confirm which execution-manager class was selected.\\
Wait returned but task continues &
Timed waits do not cancel. Distinguish a timeout, a status transition,
an async context join, and a managed goal join.\\
Goal ID disappeared &
Check history pruning and restarts. IDs are not persistent records.\\
Cancel does not stop motion &
Verify provider interruption services and blocking calls. Symbolic
cancellation is not a hardware stop implementation.\\
Schedule does not run again &
Check whether the previous child is still nonterminal, whether the
end date has passed, and whether a persistent child prevents new ticks.\\
\end{longtable}

\subsection{Focused repository checks}

The following command exercises existing language and execution
regressions without launching a robot:
\begin{lstlisting}[language=bash]
./gradlew :core:test \
  --tests '*ActionASLTest' \
  --tests '*RegressionTest' \
  --tests '*JoinOnGoalListenerTest'
\end{lstlisting}

Useful language fixtures include \code{controlTest.asl},
\code{for-loop.asl}, \code{for-each.asl}, \code{async_tests.asl},
\code{trycatch_tests.asl}, \code{multiple-returns-test.asl},
\code{observation_tests.asl}, \code{actor_dynamic_setting.asl}, and the
\code{recovery/} test directory. Some older documentation examples and
domain scripts can lag the current grammar; for a syntax question,
consult \code{ASL.g4}, its parser listener, and the relevant regression
before copying a historical example.

\subsection{Building this document}

This TeX source is self-contained and does not require the old editor
screenshot, TikZ figures, or external ASL listing files. With a LaTeX
installation providing KOMA-Script, listings, geometry, and hyperref:
\begin{lstlisting}[language=bash]
cd core/src/main/java/edu/tufts/hrilab/action/doc
pdflatex -interaction=nonstopmode -halt-on-error GoalManagerDoc.tex
pdflatex -interaction=nonstopmode -halt-on-error GoalManagerDoc.tex
\end{lstlisting}

Run twice to resolve the contents and cross-references. The TeX file is
the editable reference; an older checked-in PDF must be regenerated to
reflect changes.

\subsection{Maintaining implementation accuracy}

When extending the system, update the corresponding surface rather than
only the prose: a new ASL construct needs grammar, generated parser,
listener/writer support, a runtime context, and regression coverage.
A new public operation needs the relevant TRADE annotation and callback/
argument contract. A new action model needs coherent conditions,
effects, observation policy, and interruption behavior.

In particular, preserve the distinctions documented here: submitted
action versus desired state, model assertion versus sensed result,
selection utility versus scheduling priority, top-level goal ID versus
context ID, and acceptance of a control request versus completed physical
execution. These boundaries explain most otherwise surprising Goal
Manager behavior.

\end{document}
